% titel: Blatt 2, A6: Empirische Formel aus Verbrennungsanalyse
% skript: Mortimer Kap. 4.2 (Verbrennung, empirische Formel)

% Teil I: Werkzeuge, kein Lösungsweg.

Bei einer Verbrennung wird der Kohlenstoff der Probe vollstaendig zu
$\mathrm{CO_2}$ und der Wasserstoff zu $\mathrm{H_2O}$ umgesetzt. Aus den
Massen der Verbrennungsprodukte gewinnt man die Stoffmengen von C und H in der
Probe; ob Sauerstoff enthalten ist, folgt aus der Massenbilanz.

\begin{werkzeugkasten}[Verbrennungsanalyse]
Grundlage ist die Atomsortenerhaltung:
\begin{equation}
  n(\mathrm{C}) = n(\mathrm{CO_2})
  = \frac{m_{\mathrm{CO_2}}}{M_{\mathrm{CO_2}}},\qquad
  n(\mathrm{H}) = 2\, n(\mathrm{H_2O})
  = 2\,\frac{m_{\mathrm{H_2O}}}{M_{\mathrm{H_2O}}}.
\end{equation}
Jedes $\mathrm{H_2O}$-Molekül enthaelt zwei H. Da die Probe nur aus C, H und
möglicherweise O besteht, folgt die Masse des Sauerstoffs per Differenz:
\begin{equation}
  m_{\mathrm{O}} = m_{\mathrm{Probe}} - m_{\mathrm{C}} - m_{\mathrm{H}},\qquad
  m_i = n_i \cdot M_i,
\end{equation}
daraus $n(\mathrm{O})$. Das ganzzahlige Verhaeltnis $n(\mathrm{C}):
n(\mathrm{H}):n(\mathrm{O})$ ergibt die empirische Formel (kleinstes
ganzzahliges Verhaeltnis).
\textit{Skript:} Mortimer Kap. 4.2, Übung 2/6.
\end{werkzeugkasten}

\begin{achtung}[Darauf achten]
\begin{itemize}
  \item Der Faktor $2$ bei Wasserstoff: ein Molekül $\mathrm{H_2O}$ liefert
        zwei H-Atome, die Probe also die doppelte Stoffmenge Wasserstoff.
  \item Einheiten einheitlich waehlen (hier alle in Gramm oder alle in
        $\mathrm{mg}$); nicht $\mathrm{g}$ und $\mathrm{mg}$ mischen.
  \item Sauerstoff wird \emph{nicht} direkt gemessen, sondern aus der
        Differenz zur Probemasse bestimmt.
  \item Auf das kleinste ganzzahlige Verhaeltnis kürzen (durch die kleinste
        Stoffmenge teilen).
\end{itemize}
\end{achtung}
